Self-supervised learning (SSL) is often motivated by label scarcity: representations learned from unlabeled images are meant to make a handful of labels go further. The evidence is mostly at ImageNet scale with large networks, where contrastive and related joint-embedding methods such as SimCLR~\cite{chen2020simple}, MoCo~\cite{he2019momentum} and BYOL~\cite{grill2020bootstrap}, and pretext tasks such as rotation prediction~\cite{gidaris2018unsupervised}, are evaluated with linear probes or semi-supervised fine-tuning~\cite{zhai2019s,chen2020big}. Whether the same ordering holds for a tiny encoder and 1797 $8\times8$ images, where classical baselines such as raw pixels or a PCA projection followed by logistic regression are available, is, to our knowledge, not often reported with matched probes and several seeds.

We study exactly that question on the scikit-learn digits dataset. Our contribution is a small, controlled measurement, not a new method: (i) a SimCLR-style contrastive objective and rotation prediction, both reimplemented, are used to pretrain one shared CNN encoder on an unlabeled pool; (ii) a linear probe is trained on 10, 50 or 200 labels; (iii) results are compared with supervised scratch training of the same encoder, PCA with logistic regression, raw pixels, and a random frozen encoder, over five seeds; (iv) augmentation, temperature and pretraining-length ablations are run. We pre-specified four hypotheses (Section~\ref{sec:setup}) and report each as supported, refuted or inconclusive.
